% Track d1 -- The undivided energy and its differentiation.
% Verification: investigations/rest_energy/verify/d1_unitary_energy.py (claims d1-0 ... d1-11).
\section{The undivided energy and its differentiation}
\label{d1:sec:main}

The hypothesis reads $E=mc^2$ with $c^2$ as the symmetric part of the equation and $m$ as what breaks that symmetry; the agreed
clarification asks that energy first be one undivided, six-dimensional entity, and that mass appear through an internal differentiation
of that unity before any interaction is applied. We follow that order: the undivided energy is built from the arena alone and
differentiated with no other field and no coupling, $c$ is restored, and the Yang--Mills field $b$ \cite{YangMills54} enters only in
\cref{d1:sec:b}, where it belongs. Results already established are used without repetition: $c^2$ by itself is not a symmetry, a nonzero mass
breaks no Lorentz invariance, the unrestricted six-dimensional scalar has energy unbounded below, surplus temporal momentum lowers
four-dimensional mass, and non-orthogonality of $\bE$ and $\bB$ generates no mass (Section~14.4 of \cite{ES17}). The verdicts are collected in
\cref{d1:sec:assess}.

\paragraph{Conventions.} We use the field-strength reading of Theorem~12.3 of \cite{ES17}: $X=\sum_i(B_ie_i+E_ie_{i+3})$, $z=\bB+i\bE\in\C^3$,
the binding root $G$ ($Ge_i=e_{i+3}$, $Ge_{i+3}=-e_i$) acts as $i$, $q(z)=z\cdot z=|\bB|^2-|\bE|^2+2i\,\bE\cdot\bB$, $\eta(X,X)=\Re q$ and
$\eta(X,GX)=-2\,\bE\cdot\bB$. The frame variable is minus the textbook field, $\Btext=-\bB$ (Section~1.4 of \cite{ES17}). With
$h(z,w)=\sum_j\bar z_jw_j$ we put
\begin{equation}
\label{d1:eq:defs}
\begin{gathered}
E_6:=\tfrac12|z|^2=\tfrac12\bigl(|\bB|^2+|\bE|^2\bigr),\qquad E_0:=\tfrac12|q(z)|,\\
\mathbf S:=\bE\times\bB=\tfrac i2\,\bar z\times z\ \ \text{(frame Poynting vector, Section~12 of \cite{ES17})},\qquad
\mathbf S_{\rm t}:=\bE\times\Btext=-\mathbf S .
\end{gathered}
\end{equation}
$\mathbf S_{\rm t}$ is the physical Poynting vector and $|\mathbf S|=|\mathbf S_{\rm t}|$. Units are Heaviside--Lorentz with $c=1$ until
\cref{d1:sec:c}. For cyclic $(i,j,k)$ the diagonal rotations and boosts of $\lG$ are $L_i=\tfrac12(e_{j+3}e_{k+3}-e_je_k)$ and
$K_i=\tfrac12(e_ke_{j+3}-e_je_{k+3})$, acting on vectors by the commutator: $L_i$ by $\delta\bB=e_i\times\bB$, $\delta\bE=e_i\times\bE$, and $K_i$
by $\delta\bE=e_i\times\bB$, $\delta\bB=-e_i\times\bE$. These laws, the brackets of Theorem~12.3, Example~12.4 and Lemma~3.4 of \cite{ES17}
($\Ad_{e_1e_2e_3}|_V=\eta$, $G=-\mathcal K\circ\Ad_{e_1e_2e_3}$) were checked in the realisation of Section~2 of \cite{ES17} (claim d1-0); identifiers
d1-$n$ refer to the accompanying verification script.

\subsection{One undivided energy, with symmetry $U(3)$}
\label{d1:sec:unique}

\begin{theorem}[Uniqueness of the undivided energy]
\label{d1:thm:unique}
\begin{enumerate}
\item The symmetric forms on $V$ invariant under the diagonal rotations are spanned by $|\bB|^2$, $|\bE|^2$, $\bB\cdot\bE$, and split as
$\R\,\id\oplus\operatorname{span}\{\eta,\eta G\}$, the first summand commuting and the second anticommuting with $G$.
\item A quadratic energy invariant under rotations and under duality $e^{\vartheta G}$ is a multiple of $E_6$, which is positive definite.
\item The forms invariant under all of $\lG$ are $\operatorname{span}\{\eta,\eta G\}$, that is $|\bB|^2-|\bE|^2$ and $\bE\cdot\bB$; none is definite.
\end{enumerate}
\end{theorem}

\begin{proof}
The diagonal rotations act on $V^+\oplus V^-$ as $\R^3\otimes\R^2$ with $SO(3)$ on the first factor, which is absolutely irreducible, so by
Schur's lemma \cite{GoodmanWallach} the invariant forms are $N=\bigl(\begin{smallmatrix}\alpha&\gamma\\ \gamma&\beta\end{smallmatrix}\bigr)\otimes\id$.
With $G=\bigl(\begin{smallmatrix}0&-1\\1&0\end{smallmatrix}\bigr)\otimes\id$, $\tfrac12(N-GNG)=\tfrac12(\alpha+\beta)\id$ and
$\tfrac12(N+GNG)=\tfrac12(\alpha-\beta)\eta-\gamma\,\eta G$. Since $G^{\!\top}G=\id$ and $G^{\!\top}=-G$, $N$ is $e^{\vartheta G}$-invariant iff
$G^{\!\top}N+NG=0$ iff $[N,G]=0$, which gives (2). For (3) the boost law gives
$\delta(\alpha|\bB|^2+\beta|\bE|^2+2\gamma\,\bB\cdot\bE)=-2(\alpha+\beta)\,\bv\cdot(\bE\times\bB)$, the $\gamma$ term being unchanged, so $\beta=-\alpha$.
\end{proof}

This is Theorem~17.11(iii) of \cite{ES17} for forms on $V$, with the Lorentz class of Proposition~17.4. It is fixed: it could have failed, since
an $\eta$-self-adjoint complex structure need not be orthogonal for a positive form (conjugating $G$ by a boost outside $\lG$ destroys this), and
then no positive duality-invariant energy would be rotation invariant (d1-1a). The rotation-invariant forms carry duality weights $0$ and
$\pm2$: $E_6$ has weight $0$, while $e^{\vartheta G}$ turns $(\Re q,\Im q)$ through $2\vartheta$. In this precise sense $E_6$ is undivided: it does
not distinguish the magnetic triple from the electric one, nor, as we now show, rest from motion.

\begin{proposition}[The symmetry of the undivided energy is $U(3)$]
\label{d1:prop:u3}
(1) The maps of $V$ preserving $E_6$ and commuting with $G$ form $\uu3$, of dimension $9$: $A=\bigl(\begin{smallmatrix}P&Q\\-Q&P\end{smallmatrix}\bigr)$
acts on $z$ by the anti-Hermitian matrix $P-iQ$. (2) The Lie algebra of $O(3,\C)\times U(1)_G$, the $G$-linear maps preserving $q$ up to a phase, is
$\lG\oplus\R G$, of dimension $7$. (3) Their intersection is $\so3\oplus\R G$, of dimension $4$, the maximal compact subalgebra of $\lG\oplus\R G$; and
$\uu3\cap\lG=\so3$. (4) The stabiliser of the function $E_0$ in $U(3)$ is $U(1)\cdot O(3)=\{e^{i\varphi}O\}$; for instance the unitary map
$(1,0,0)\mapsto\tfrac1{\sqrt2}(1,i,0)$ carries a field at rest ($E_0=E_6=\tfrac12$) to a null field ($E_0=0$) at the same $E_6$.
\end{proposition}

\begin{proof}
(1) Commuting with $G$ forces the block form, which is multiplication by $P-iQ$ on $z$; $A^{\!\top}=-A$ means $P^{\!\top}=-P$, $Q^{\!\top}=Q$, so $P-iQ$
is anti-Hermitian; $3+6=9$. (2) $q(e^{tM}z)=e^{i\varphi t}q(z)$ gives $M^{\!\top}+M=i\varphi\,\id$, so $M=\tfrac{i\varphi}2+M_0$ with $M_0\in\so{3,\C}=\lG$
(Theorem~12.1 of \cite{ES17}), and $i\,\id$ is $G$. (3) A complex antisymmetric anti-Hermitian matrix is real. (4) If $U\in U(3)$ preserves $|q|$, the
form $z^{\!\top}U^{\!\top}Uz$ vanishes on the cone $q=0$; $q$ is irreducible, so $U^{\!\top}U=e^{i\varphi}\id$, and $O=e^{-i\varphi/2}U$ is unitary and complex
orthogonal, hence real.
\end{proof}

Parts (1)--(3) are inherited Hermitian linear algebra (the arena contributes the orthogonality of $G$), and they make ``unitary'' literally true:
the symmetry group of the undivided energy, compatible with the binding root, is $U(3)$. Part (4) is fixed and is the exact sense in which the rest
energy ``breaks'' something: $U(3)$ (dimension $9$) is reduced to $U(1)\cdot O(3)$ (dimension $4$), and the five-dimensional orbit
$U(3)/U(1)\cdot O(3)$ parametrises the real structures $z\mapsto\bar z$ up to phase, of which the arena has already chosen one
(\cref{d1:prop:Tparity}). Verified exactly (d1-1b, d1-8b).

\subsection{The rest energy is the complete invariant of Lorentz and duality}
\label{d1:sec:orbits}

\begin{theorem}[Orbits]
\label{d1:thm:orbits}
(1) Each level set $\{|q|=r^2\}$, $r>0$, is one orbit of $O(3,\C)\times U(1)_G$, of real dimension $5$, with two-dimensional stabilisers; every
non-null $z$ is carried to $(r,0,0)$, $r=\sqrt{2E_0}$. (2) The nonzero null fields form one orbit of $SO(3,\C)$ alone, of dimension $4$, and at a null
point $Gz$ is tangent to the Lorentz orbit. (3) Hence $E_0$ is a complete invariant of Lorentz and duality on non-null fields, and $E_6$ is the
complete invariant of $U(3)$, whose orbits are the spheres.
\end{theorem}

\begin{proof}
(1) Write $q(z)=r^2e^{i\varphi}$ and $w=e^{-i\varphi/2}z/r$, so $q(w)=1$; the factor is the duality rotation $e^{-\varphi G/2}$. As $q(w)\neq0$, $q$ is
nondegenerate on $w^\perp=\{f:f\cdot w=0\}$, so some $f\in w^\perp$ has $f\cdot f\neq0$. Put $f_2=f/\sqrt{f\cdot f}$ and $f_3=w\times f_2$; the bilinear
identity $(a\times b)\cdot(a\times b)=(a\cdot a)(b\cdot b)-(a\cdot b)^2$ gives $f_3\cdot f_3=1$. Then $M=(w\,f_2\,f_3)$ has $M^{\!\top}M=\id$ and
$\det M=(w\times f_2)\cdot f_3=1$, and $A=M^{\!\top}\in SO(3,\C)$ satisfies $Ae^{-i\varphi/2}z=re_1$. An element $M_0+\tfrac{i\varphi}2$ fixes $e_1$ iff
$\varphi=0$ and $M_0$ is a complex rotation about $e_1$, since $M_0e_1\in\operatorname{span}_\C\{e_2,e_3\}$: the stabiliser is $\so{2,\C}$ and the orbit
has dimension $7-2=5$. (2) A null $z=\mathbf a+i\mathbf b$ has $|\mathbf a|=|\mathbf b|>0$, $\mathbf a\cdot\mathbf b=0$; a diagonal rotation carries it to
$|\mathbf a|(1,i,0)$. The boost $K_3$ acts by $\delta z=i\,e_3\times z$ and $i\,e_3\times(1,i,0)=(1,i,0)$, so $e^{\chi K_3}$ with $\chi=-\log|\mathbf a|$ (the
Doppler factor) finishes. The Lorentz tangent space $\{\mathbf a\times z\}$ at a null $z$ is its $q$-orthogonal complement, which contains $z$ and $iz$;
at $(1,i,0)$, $iz=-e_3\times z$ is the rotation $L_3$, the helicity of Example~13.11 of \cite{ES17}. (3) follows.
\end{proof}

Inherited: the classification of a field strength at a point \cite{Synge56,LandauLifshitz2}, resting on Theorem~12.1 of \cite{ES17}; checked
constructively at $2000$ random points and by exact tangent ranks (d1-2a, d1-2b). $U(3)$ sees only $E_6$, Lorentz with duality only $E_0$, and the two
share the maximal compact $U(1)\cdot O(3)$; without duality the complete invariant is $q$, whose phase (twice the complexion of Proposition~12.12 of
\cite{ES17}) duality removes.

\subsection{The differentiation identity}
\label{d1:sec:identity}

\begin{proposition}[Differentiation]
\label{d1:prop:lagrange}
For every $z\in\C^3$
\begin{equation}
\label{d1:eq:split}
|z|^4=|z\cdot z|^2+|\bar z\times z|^2,\qquad\text{equivalently}\qquad E_6^2=E_0^2+|\mathbf S|^2 .
\end{equation}
The rank-one Hermitian $H=zz^\dagger$ carries all three terms: $\tr H=2E_6$, $(\Im H)_{jk}=\epsilon_{jkl}S_l$, $\tr(H\bar H)=|q|^2=4E_0^2$, $H^2=(\tr H)H$.
Hence $E_0\le E_6$, with equality iff $\mathbf S=0$, iff $\bar z\parallel z$, iff $\bB\parallel\bE$ or one of them vanishes.
\end{proposition}

\begin{proof}
$\bar z\times z=-2i\,\bE\times\bB$. The identity $(a\times b)\cdot(c\times d)=(a\cdot c)(b\cdot d)-(a\cdot d)(b\cdot c)$ at $(\bar z,z,z,\bar z)$ gives
$(\bar z\times z)\cdot(z\times\bar z)=|z|^4-|q|^2$, and $z\times\bar z$ is the conjugate of $\bar z\times z$. Next
$H_{jk}=B_jB_k+E_jE_k+i(E_jB_k-B_jE_k)$ with $E_jB_k-B_jE_k=\epsilon_{jkl}(\bE\times\bB)_l$, and $\tr(H\bar H)=|\sum_jz_j^2|^2$. The equality case is
Cauchy--Schwarz for $|h(\bar z,z)|\le|z|^2$.
\end{proof}

Inherited, a complex Lagrange identity (d1-3a). In the user's terms it is exact: the undivided $E_6$ is the hypotenuse of a hyperbolic triangle
whose legs are a Lorentz-invariant rest part $E_0$ and a motion part $|\mathbf S|$, and the split needs no interaction and no second field.

\begin{proposition}[Rainich identity]
\label{d1:prop:rainich}
For the stress--energy tensor of $(\bE,\Btext)$, $T^{00}=\tfrac12(|\bE|^2+|\Btext|^2)=E_6$, $T^{0i}=(\mathbf S_{\rm t})_i$,
$T^{ij}=T^{00}\delta_{ij}-E_iE_j-B_{{\rm t}i}B_{{\rm t}j}$, one has $T^\mu{}_\mu=0$ and
\begin{equation}
\label{d1:eq:rainich}
T^\mu{}_\lambda T^\lambda{}_\nu=E_0^2\,\delta^\mu{}_\nu ,\qquad T_{\mu\nu}T^{\mu\nu}=4E_0^2 .
\end{equation}
$T^{\mu\nu}$ is duality invariant, and for every unit timelike observer $t$ the density $P(t)^\mu=T^\mu{}_\nu t^\nu$ is future causal with
$\eta(P(t),P(t))=E_0^2$.
\end{proposition}

\begin{proof}
Duality invariance is a polynomial identity, and both sides of \eqref{d1:eq:rainich} are Lorentz covariant, so by \cref{d1:thm:orbits} it suffices to take
$z=(r,0,0)$, where $T^\mu{}_\nu=\tfrac{r^2}2\operatorname{diag}(1,1,-1,-1)$ and $E_0=\tfrac{r^2}2$; null fields follow by continuity. Then
$\eta(Tt,Tt)=E_0^2\eta(t,t)$, and causality of $P(t)$ is the dominant energy condition \cite{HawkingEllis73}.
\end{proof}

Inherited \cite{Rainich1925,Synge56}; verified symbolically (d1-3b). The mass of a field point is the Lorentz scalar of the same tensor whose $00$
component is the undivided energy. But $(u,\mathbf S_{\rm t})$ is not a four-vector: each observer sees a different density $P(t)$, all of length $E_0$.

\subsection{The rest energy as a minimum, and two boost factors}
\label{d1:sec:min}

\begin{theorem}[Variational form of the rest energy]
\label{d1:thm:min}
\begin{enumerate}
\item For non-null $z$, $E_0$ is the minimum of $E_6$ over the Lorentz orbit, and over the Lorentz--duality orbit, attained exactly where
$\mathbf S=0$; for null $z\neq0$ the infimum $0$ is not attained. Moreover $E_6=E_0/\sqrt{1-v_E^2}$ with $v_E:=|\mathbf S|/E_6<1$.
\item If $\bB\parallel\bE\parallel\mathbf n$ and the boost $e^{\chi K}$ of rapidity $\chi$ has axis orthogonal to $\mathbf n$, then $E_6=E_0\cosh2\chi$,
$|\mathbf S|=E_0\sinh2\chi$ and $v_E=\tanh2\chi=2\beta/(1+\beta^2)$, $\beta=\tanh\chi$; a boost along $\mathbf n$ changes neither $E_6$ nor $\mathbf S=0$.
\item The observer moving with $\mathbf v_{\rm rest}=\mathbf S_{\rm t}/(E_6+E_0)$ sees $\mathbf S=0$ and $E_6=E_0$; $|\mathbf v_{\rm rest}|=\beta$.
\end{enumerate}
\end{theorem}

\begin{proof}
(1) By \eqref{d1:eq:split} $E_6\ge E_0$ with equality iff $\mathbf S=0$, and \cref{d1:thm:orbits} supplies points with $\mathbf S=0$ in both orbits; for
null $z$ the Doppler boost drives $E_6$ to $0$. Divide \eqref{d1:eq:split} by $E_6^2$. (2) Take $\mathbf n=e_1$, $X=ae_1+be_4$. By Example~12.4 of
\cite{ES17} and the boost law $[K_3,e_1]=e_5$, $[K_3,e_5]=e_1$, $[K_3,e_4]=-e_2$, $[K_3,e_2]=-e_4$, so $\ad_{K_3}^3=\ad_{K_3}$ and
$e^{\chi\ad_{K_3}}=\id+\sinh\chi\,\ad_{K_3}+(\cosh\chi-1)\ad_{K_3}^2$. Hence $\bB=(a\cosh\chi,-b\sinh\chi,0)$, $\bE=(b\cosh\chi,a\sinh\chi,0)$,
$E_6=\tfrac12(a^2+b^2)\cosh2\chi$, $\mathbf S=-\tfrac12(a^2+b^2)\sinh2\chi\,e_3$, $E_0=\tfrac12(a^2+b^2)$; likewise for $K_2$, while $K_1$ fixes $e_1$ and
$e_4$. (3) $|\mathbf S|/(E_6+E_0)=\sinh2\chi/(\cosh2\chi+1)=\tanh\chi$. For the direction: $e^{\chi K_3}$ is the passage to a frame moving with
$-\tanh\chi\,e_3$ (Section~12 of \cite{ES17}), in which the rest frame moves along $+e_3$, and so does $\mathbf S_{\rm t}=-\mathbf S$.
\end{proof}

Inherited \cite{LandauLifshitz2}; verified with the paper's own $K_i$ symbolically, by minimisation at $25$ random fields, and at $100$ random fields
for (3), where the boost with $\mathbf S_{\rm t}/(E_6+E_0)$ leaves a Poynting vector below $10^{-13}E_0$ and the one built from $\bE\times\bB$ leaves at
least $0.2\,E_0$ (d1-4a--d1-4c). So the rest-frame velocity $\mathbf S/(u+\tfrac12|q|)$ and the relation $\tfrac12|q|=u\sqrt{1-v_E^2}$ of
Section~14.4 of \cite{ES17} are correct provided $\mathbf S$ is the textbook $\bE\times\Btext$; with the frame vector $\bE\times\bB$, which Section~12
also calls the Poynting vector, the direction is reversed. And a field point carries two boost factors, $E_6=E_0\cosh2\chi$ against $E=mc^2\cosh\chi$
for a particle, because $T^{\mu\nu}$ has rank two: the rest frame moves with $\beta$, the energy with $\beta\oplus\beta$, and an observer riding with
the energy velocity sees the same $E_6$ and the reversed Poynting vector (d1-4b). $E_0$ is a rest-energy density of a field point, not the rest energy
of a particle carried by the flux.

\subsection{Restoring $c$: where $c^2$ enters}
\label{d1:sec:c}

Adding the dimensionless $G$ to a derivative along a clock coordinate needs a length $\ell$ (Section~3 of \cite{ES17}):
\begin{equation}
\label{d1:eq:kappa}
\kappa=c/\ell,\qquad \alpha_\pm=\partial_\tau\pm\kappa G .
\end{equation}

\begin{proposition}[$c^2$ is an isotropic weight]
\label{d1:prop:onec}
Let the temporal triple carry $\bB$ and the spatial triple $\bE$, $C:V^+\to V^-$ the conversion and
$G_C=\bigl(\begin{smallmatrix}0&-C^{-1}\\ C&0\end{smallmatrix}\bigr)$, $G_C^2=-\id$. (1) $G_C$ commutes with the diagonal rotations iff $C=c\,\id$.
(2) For $C=c\,\id$ the forms invariant under the diagonal rotations and under $e^{\vartheta G_c}$ are the multiples of $c^2|\bB|^2+|\bE|^2$.
(3) $(\kappa G)^2=-c^2/\ell^2$, and $\ell$ is not fixed.
\end{proposition}

\begin{proof}
(1) The lower-left block of $[G_C,R\oplus R]$ is $CR-RC$, and Schur's lemma gives $C=c\,\id$: Theorem~3.2 of \cite{ES17}. (2) For
$N=\bigl(\begin{smallmatrix}\alpha&\gamma\\ \gamma&\beta\end{smallmatrix}\bigr)\otimes\id$,
$G_c^{\!\top}N+NG_c=\bigl(\begin{smallmatrix}2c\gamma&c\beta-\alpha/c\\ c\beta-\alpha/c&-2\gamma/c\end{smallmatrix}\bigr)\otimes\id$, zero iff
$\gamma=0$ and $\alpha=c^2\beta$. (3) $G^2=-\id$.
\end{proof}

Fixed by Theorem~3.2 of \cite{ES17}, and verified exactly, including the failure for $C=\operatorname{diag}(c_1,c_2,c_3)$ with distinct entries (d1-5b).

\begin{proposition}[SI form]
\label{d1:prop:SI}
With $\varepsilon_0\mu_0c^2=1$, $z_c:=c\bB+i\bE=-c\Btext+i\bE$, $u=\tfrac{\varepsilon_0}2(|\bE|^2+c^2|\Btext|^2)=\tfrac{\varepsilon_0}2|z_c|^2$,
$\mathbf g=\mathbf S_{\rm t}/c^2$ with $\mathbf S_{\rm t}=\bE\times\Btext/\mu_0$, $q_c:=z_c\cdot z_c=c^2|\Btext|^2-|\bE|^2-2ic\,\Btext\cdot\bE$
$(=c^2|\bB|^2-|\bE|^2+2ic\,\bB\cdot\bE)$ and $u_0:=\tfrac{\varepsilon_0}2|q_c|$,
\begin{equation}
\label{d1:eq:SI}
u^2=c^2|\mathbf g|^2+u_0^2,\qquad \rho_0:=u_0/c^2,
\end{equation}
with $[u]=[c\,\mathbf g]=[u_0]=\mathrm{J\,m^{-3}}$, $[\rho_0]=\mathrm{kg\,m^{-3}}$, $[q_c]=\mathrm{V^2\,m^{-2}}$.
\end{proposition}

\begin{proof}
$c^2|\mathbf g|^2=\varepsilon_0^2c^2(|\bE|^2|\Btext|^2-(\bE\cdot\Btext)^2)$ and $\tfrac{\varepsilon_0^2}4|q_c|^2=\tfrac{\varepsilon_0^2}4[(c^2|\Btext|^2-|\bE|^2)^2+4c^2(\bE\cdot\Btext)^2]$
add up to $\tfrac{\varepsilon_0^2}4(|\bE|^2+c^2|\Btext|^2)^2$.
\end{proof}

Inherited (d1-5a). This is the exact content of ``$c^2$ is the symmetric part'': $c$ is the unique conversion between the triples that commutes with
their common rotations, the residual symmetry of both $E_6$ and $E_0$, and it enters the undivided energy only as the isotropic weight $c^2$ of
$|\bB|^2$. It is a unit conversion, not a symmetry, and nothing breaks it: $u_0$ uses the same $c$, and $\rho_0=u_0/c^2$ only converts units.

\subsection{Finite configurations: the mass of a whole}
\label{d1:sec:whole}

\begin{proposition}[Mass of a whole]
\label{d1:prop:whole}
\begin{enumerate}
\item For a free field with compactly supported Cauchy data, $(U,\mathbf P)=\int(u,\mathbf S_{\rm t})\,d^3x$ is a four-vector and
$M^2:=U^2-|\mathbf P|^2\ge U\int(u-|\mathbf S|)\,d^3x\ge0$, so $M>0$ whenever $q\neq0$ on a set of positive measure.
\item In every frame $M\ge\int E_0\,d^3x$, with equality iff the densities $(u,\mathbf S_{\rm t})(x)$ are all proportional.
\item A plane wave has $M=0$. The waves $E_x=a\cos(kz\mp\omega t)$, $cB_{{\rm t}y}=\pm a\cos(kz\mp\omega t)$ superpose to $E_x=2a\cos kz\cos\omega t$,
$cB_{{\rm t}y}=2a\sin kz\sin\omega t$; per cell (unit area, length $2\pi/k$) $U=2\pi\varepsilon_0a^2/k$ and $\mathbf P=0$, so $Mc^2=U$, although each
constituent has $q_c=0$ and $U_\pm=|\mathbf P_\pm|c$. The total has $q_c=-2a^2(\cos2kz+\cos2\omega t)$, $\int_{\rm cell}T^{zz}=U$,
$\int_{\rm cell}T^{xx}=-U\cos2\omega t$, and the integrated rest energy $\int_{\rm cell}u_0$ oscillates between $\tfrac2\pi U$ and $U$.
\item With $c=1$ and $f(s)=(1-s^2)^2$ on $|s|\le1$, $0$ elsewhere, $E_x=f(z-t)+f(z+t)$, $B_{{\rm t}y}=f(z-t)-f(z+t)$ has $U=\tfrac{512}{315}$ per unit
area and $\mathbf P=0$ at all times, so $M=U$; for $|t|\ge1$ the pulses are disjoint, $q\equiv0$ and $\int E_0=0$.
\end{enumerate}
\end{proposition}

\begin{proof}
(1) Conservation and finite propagation speed make $\int_\Sigma T^{\mu\nu}d\Sigma_\nu$ independent of the spacelike hyperplane $\Sigma$
\cite{Laue1911}; $u-|\mathbf S|=E_0^2/(u+|\mathbf S|)\ge0$, and $M^2=(U-|\mathbf P|)(U+|\mathbf P|)$. (2) Each $(u,\mathbf S_{\rm t})(x)$ is future causal of
length $E_0(x)$ (\cref{d1:prop:rainich}), and future causal vectors obey $|P_1+P_2|\ge|P_1|+|P_2|$, with equality iff proportional. (3) Elementary
cell integrals; for the last statement, with $s=\sin^2\omega t$ one has $q_c=4a^2(s-\cos^2kz)$, so the integrand is $u_0=2\varepsilon_0a^2|s-\cos^2kz|$, whose mean is least at the median $s=\tfrac12$ of the
arcsine law, giving $\tfrac2\pi U$. (4) $u=f(z-t)^2+f(z+t)^2$ and $S_z=f(z-t)^2-f(z+t)^2$, the cross terms cancelling, and $q=-4f(z-t)f(z+t)$.
\end{proof}

Items (1) and (3) are inherited, (1) being the statement of Section~14.4 of \cite{ES17}. The cell in (3) is not isolated: $\int T^{zz}=U$ is the
radiation pressure on whatever confines it, so its $(U,\mathbf P)$ is a four-vector only with the walls included \cite{Laue1911}, as in Einstein's box
of light \cite{Einstein1905}. Items (2) and (4) refute the dictionary's reading that the mass of a whole ``emerges from the interference that makes
$q\neq0$'': interference suffices for $M>0$ but is not necessary, since disjoint counter-propagating null pulses have $q\equiv0$ and $M=U>0$. A whole
gets its mass from the non-collinearity of its constituents' energy--momenta, as two photons do, and $M$ dominates the integrated rest energy of its
points; the interference reading is obstructed. $M$ is also blind to orthogonality: the data $(\bE,\Btext)=(0,\mathbf b)$ and
$(\mathbf b\sin\vartheta,\mathbf b\cos\vartheta)$ have the same $T^{\mu\nu}$ while $\bE\cdot\Btext$ is $0$ and $\tfrac12\sin2\vartheta|\mathbf b|^2$
(d1-6a--d1-6c).

\subsection{The momentum reading of ``$E$ is six-dimensional''}
\label{d1:sec:momentum}

\begin{proposition}[Splitting the temporal energy lowers the mass]
\label{d1:prop:momentum}
Let $P=\sum_i(E_i/c)e_i+\sum_ap_ae_{a+3}$ with $\eta(P,P)=\mu^2c^2$, $\widehat W$ the symmetric clock of Proposition~3.5 of \cite{ES17},
$E_W:=(E_1+E_2+E_3)/\sqrt3$ and $\mathbf E_\perp$ the part of $(E_1,E_2,E_3)$ orthogonal to $(1,1,1)$, the relative-time energies. The mass read by the
clock, $m_4^2c^4:=E_W^2-|\bp|^2c^2$, is
\[
m_4^2c^4=\mu^2c^4-|\mathbf E_\perp|^2\le\mu^2c^4 ;
\]
for the clock $e_1$, $m_4^2c^4=\mu^2c^4-E_2^2-E_3^2$, and $E=(3,4,0)$, $\bp=(5,0,0)$, $c=1$ give $\mu=0$, $m_4^2=-16$.
\end{proposition}

\begin{proof}
$|(E_1,E_2,E_3)|^2=E_W^2+|\mathbf E_\perp|^2$ and $\mu^2c^4=c^2\eta(P,P)=|(E_1,E_2,E_3)|^2-|\bp|^2c^2$.
\end{proof}

Obstructed, the clock-covariant form of Proposition~27.1 and Example~27.3 of \cite{ES17} (d1-7). In the momentum reading the differentiation of the
temporal energy subtracts from the mass and can make it tachyonic, and under the surplus constraints of Theorem~30.1 of \cite{ES17},
$\mathbf E_\perp=0$, nothing is left to differentiate; in the field-strength reading the frame-dependent part adds, $E_6\ge E_0$, as momentum should.
Only the field-strength reading of ``$E$ is six-dimensional'' carries the structure the user describes.

\subsection{Three candidates for $m$ in the arena's own structures}
\label{d1:sec:m}

\begin{proposition}[The mass form acts on the frame as parity]
\label{d1:prop:Tparity}
Let $T=e_1e_2e_3$ and $P:=\Ad_T|_V$. (1) $P=\eta$ fixes $\bB$ and reverses $\bE$: it is parity on the frame; it commutes with the $L_i$ and
anticommutes with $G$ and the $K_i$, so it normalises $\lG$ without being $\lG$-invariant. (2) $PX$ corresponds to $\bar z$ and $q(z)=h(Pz,z)$, so
$E_6=\tfrac12h(z,z)$ and $E_0=\tfrac12|h(Pz,z)|=\tfrac12\max_\vartheta\eta(e^{\vartheta G}X,e^{\vartheta G}X)$. (3) As a Grassmann mass kernel $T$ is
invariant under all of $\so{3,3}$ in the sense of Definition~18.2 of \cite{ES17} and breaks nothing; only its sign changes, under conjugation by a
spatial generator.
\end{proposition}

\begin{proof}
(1) Lemma~3.4 of \cite{ES17}; $\eta G\eta=-G$ blockwise, the $L_i$ preserve each triple and the $K_i$ exchange them. (2) $z(\eta X)=\bB-i\bE$,
$h(\bar z,z)=\sum_jz_j^2$, and $\eta(e^{\vartheta G}X,e^{\vartheta G}X)=\Re(e^{2i\vartheta}q)$. (3) By Lemma~2.1 of \cite{ES17}, $e_ae_b$ commutes with
$e_1e_2e_3$ when $a,b$ have the same sign and anticommutes otherwise, and $T^{\!\top}=-T$; this is Corollary~18.4 of \cite{ES17}.
\end{proof}

The identities are fixed (d1-8a); their consequence for the dictionary is negative. ``$m\leftrightarrow T$'' fails on the frame, where $T$ is parity;
the Dirac mass lives on the interior reading, which Proposition~12.6 of \cite{ES17} excludes on the same six directions, and even there the mass term
of the full algebra breaks nothing. What survives is a reading: the rest energy is the overlap of a field with its mirror image, zero for a null field
and maximal for a field at rest; and $P$, the arena's splitting into triples, is what reduces $U(3)$ to $U(1)\cdot O(3)$ (\cref{d1:prop:u3}(4)).

\begin{proposition}[The interior energy is not a rest energy]
\label{d1:prop:interior}
For the split-form interior with the binding root (Theorem~16.2 of \cite{ES17}) every solution is $z(\tau)=c_-e^{-i\kappa\tau}+c_+e^{i\kappa\tau}$, and
$\ddot X=-\kappa^2X$ does not involve $m$. The Noether energy $\mathcal E=2m\kappa^2\Re(c_-\cdot c_+)$ is indefinite and vanishes identically on the
orbits without momentum, $c_+=0$, the duality rotations of Proposition~16.5, along which $E_6=\tfrac12|c_-|^2$ and $E_0=\tfrac12|c_-\cdot c_-|$ are
constant; in general $|z|^2=|c_-|^2+|c_+|^2+2\Re\bigl(h(c_-,c_+)e^{2i\kappa\tau}\bigr)$.
\end{proposition}

\begin{proof}
With $\eta(X,X)=\Re q$ and $\eta(GX,GX)=-\eta(X,X)$, $\mathcal E=\tfrac m2\Re(\dot z\cdot\dot z+\kappa^2z\cdot z)$ and
$\dot z\cdot\dot z+\kappa^2z\cdot z=4\kappa^2c_-\cdot c_+$; the momentum corresponds to $\dot z+i\kappa z=2i\kappa c_+e^{i\kappa\tau}$.
\end{proof}

Obstructed as an identification, verified symbolically (d1-9): in the interior $m$ is an overall factor, the purely internal motion carries zero energy,
and $c^2$ appears only in $(\kappa G)^2=-c^2/\ell^2$ with a free length. The entry ``$c^2\leftrightarrow\kappa G$'' is right that $\kappa G$ lies in the
factorising summand of Theorem~17.3 of \cite{ES17}, but nothing in the interior turns $\kappa^2$ into a rest energy. The closure does better.

\begin{theorem}[In the closure the rest energy breaks the duality translation]
\label{d1:thm:closure}
In the Lie--Poisson closure of Proposition~17.8 of \cite{ES17}, $H=h^2/(2I_0)+\tfrac12p^{\!\top}Np$, with momentum field $\pi=\eta^{-1}p$: (1) for
every $N$ the Casimir $|p|^2$ equals $2E_6(\pi)$, and $E_0(\pi)=\tfrac12|w(p)|$ and $\mathbf S(\pi)$ are constants of motion; (2) for Lorentz-class
$N=(a\eta^{-1}+bG\eta^{-1})/m$ the reduced Hamiltonian is
\[
H=\frac{h^2}{2I_0}+\frac{|a+ib|\,E_0(\pi_0)}m\cos(2\theta+\varphi_0),
\]
so the rotor's translation $\theta\mapsto\theta+{\rm const}$, the duality rotation of the momentum, is broken with amplitude $|a+ib|E_0(\pi_0)/m$, and
the rate $\kappa=h/I_0$ is constant iff $E_0(\pi_0)=0$; (3) for the duality-class $N=\id/m$ the frame energy is $E_6(\pi)/m$ and the rate is conserved
for all data.
\end{theorem}

\begin{proof}
Along the flow $p=e^{-\theta G}p_0$, so $\pi=e^{\theta G}\eta p_0$ and $z(\pi)=e^{i\theta}z(\pi_0)$, a duality rotation, which preserves $E_6$, $E_0$, $\mathbf S$;
and $|p|=|\pi|$ as $\eta$ is orthogonal. (2) is Theorem~17.13 of \cite{ES17} with $R_0=|a+ib|\,|w(p_0)|=2|a+ib|E_0(\pi_0)$, the last clause Corollary~17.14.
(3) $N=\id/m$ commutes with $G$, so $\dot h=-p^{\!\top}GNp=0$.
\end{proof}

Fixed, as a consequence of Theorem~17.13 and Corollary~17.14 of \cite{ES17}; it could have failed, and does for the duality class. Verified by
integrating the full system at four random momenta (relative conservation errors below $10^{-9}$; rate constant to $10^{-10}$ on null data and for
$N=\id/m$; d1-10). Here the user's sentence becomes a theorem: reading the clamped rate $\kappa=c/\ell$ as the ``$c$'' of the interior, that rate keeps
its symmetry (constancy, invariance under the rotor's duality translation) exactly when the momentum field is null, the rest energy $E_0(\pi)$ is the
amplitude that breaks it, and the undivided $E_6(\pi)$ is a Casimir and breaks nothing. But the breaking is explicit, from an inertia that
Theorem~17.11 of \cite{ES17} forbids from also conserving the rate; it breaks duality, not Lorentz invariance; and with the undivided energy as
Hamiltonian ($N=\id/m$) the rate is constant at the price of the frame drift of Theorem~17.12.

\subsection{Where the $b$-field belongs}
\label{d1:sec:b}

\begin{proposition}[Transport preserves the totals]
\label{d1:prop:bfield}
(1) $A\in\so{3,3}$ preserves $\bE\cdot\bB$ for all $X$, equivalently $q$, equivalently $E_0$, iff $[A,G]=0$, that is $A\in\lG$. (2) At three
constituents the binding-compatible $\so{4}_\Delta\cong\su{2}\oplus\su{2}$ of Theorem~17.21 of \cite{ES17} is compact and lies in
$\Cent(G_3)\cong\so{9,\C}$, so it preserves the total undivided energy $\tfrac12|Z|^2=\tfrac12\sum_k|z^{(k)}|^2$ and the total invariant
$q_3(Z)=\sum_kz^{(k)}\cdot z^{(k)}$. (3) It does not preserve the separate constituents' invariants: the rotation through $\pi/4$ generated by
$\tfrac12\mathrm{rot}(7,13)$ maps $(z^{(2)},z^{(3)})=((1,i,0),(1,0,0))$ to $((0,i,0),(\sqrt2,0,0))$, keeping $\sum_kE_6^{(k)}=\tfrac32$ and $q_3=1$, while
$\sum_kE_0^{(k)}$ goes from $\tfrac12$ to $\tfrac32$.
\end{proposition}

\begin{proof}
(1) For $A\in\so{3,3}$, $A^{\!\top}\eta G+\eta GA=\eta[G,A]$. If $A$ preserves $\Re q$ and $|q|$, then $\Im q\cdot\delta(\Im q)=0$, so $\delta(\Im q)=0$ since
polynomials form an integral domain. (2) Theorems~17.19 and 17.21 of \cite{ES17}; compact elements of $\so{9,\C}$ are real antisymmetric. (3) Direct,
with $[e_Ae_B,e_C]=2\eta_{BC}e_A-2\eta_{AC}e_B$, checked against the realisation.
\end{proof}

Verified exactly (d1-11); its content is an obstruction: no transport valued in the compact part of $\Cent(G_n)$, the binding-compatible $b$-field
included, creates or destroys total rest energy; it moves rest energy between constituents, and can turn one constituent's motion into another's rest
energy at fixed totals. That is where the Yang--Mills field belongs here: redistribution among constituents, once the structure of Theorem~17.21 is in
place, not the undivided energy of one constituent, whose differentiation needed no connection. On one constituent a connection changes $E_0$ only
through its component outside $\lG$, which is not a Lorentz transformation of the field.

\subsection{Assessment}
\label{d1:sec:assess}

\begin{table}[t]
\centering
\small
\begin{tabular}{@{}lll@{}}
\toprule
claim & statement & class\\
\midrule
d1-1a & $E_6$ is the unique rotation- and duality-invariant energy & fixed\\
d1-1b & its symmetry is $U(3)$; $\uu3\cap(\lG\oplus\R G)=\so3\oplus\R G$ & inherited\\
d1-2a,b & $E_0$ complete invariant of Lorentz $\times$ duality, $E_6$ of $U(3)$ & inherited\\
d1-3a,b & $E_6^2=E_0^2+|\mathbf S|^2$; Rainich identity & inherited\\
d1-4a--c & $E_0=\min E_6$; $E_6=E_0\cosh2\chi$; rest-frame velocity with $\mathbf S_{\rm t}$ & inherited\\
d1-5a & $u^2=c^2|\mathbf g|^2+u_0^2$ & inherited\\
d1-5b & $c$ enters only as the isotropic weight $c^2$ (Theorem~3.2) & fixed\\
d1-6a,c & standing wave $Mc^2=U$; $M$ blind to orthogonality & inherited\\
d1-6b & the mass of a whole is not interference; $M\ge\int E_0$ & obstructed\\
d1-7 & the momentum reading lowers the mass & obstructed\\
d1-8a & $T$ acts on the frame as parity; ``$m\leftrightarrow T$'' fails there & obstructed\\
d1-8b & $\Stab_{U(3)}(E_0)=U(1)\cdot O(3)$ & fixed\\
d1-9 & the interior Noether energy is not a rest energy & obstructed\\
d1-10 & closure: $E_0(\pi)$ is the duality-breaking amplitude & fixed\\
d1-11 & binding-compatible $b$ redistributes, cannot create, rest energy & obstructed\\
d1-8c & ``$E$ undivided, $m$ differentiated, $c^2$ symmetric'' & interpretive\\
\bottomrule
\end{tabular}
\caption{The claims of \cref{d1:sec:main} and their classes; identifiers refer to the verification script.}
\label{d1:tab:claims}
\end{table}

Most of what was asked is inherited (\cref{d1:tab:claims}), ordinary electrodynamics in the arena's variables, which by Corollary~4.4 of
\cite{ES17} predicts nothing. The arena fixes four things that could have failed: the orthogonality of its binding root, hence a unique undivided
energy; a single conversion entering as $c^2|\bB|^2$; the stabiliser $U(1)\cdot O(3)$ of the rest energy in $U(3)$, cut out by its own splitting
$P=\Ad_{e_1e_2e_3}$; and, in the closure, the rest energy of the momentum as the exact amplitude of duality breaking. Five readings are obstructed.

Corrected so that it is true, the user's sentence becomes the following interpretive statement (d1-8c), resting on the results above. The undivided
energy is $E_6=\tfrac12|z|^2$, the one quadratic energy of a six-component frame vector invariant under rotations and duality, with symmetry $U(3)$. It
differentiates internally, with no interaction, as $E_6^2=E_0^2+|\mathbf S|^2$, into a Lorentz-invariant rest energy $E_0=\tfrac12|z\cdot z|$ and a
motion part $|\mathbf S|$. What differentiates is the splitting of the six directions into a temporal and a spatial triple, the real structure
$\Ad_{e_1e_2e_3}$; what it breaks is the compact $U(3)$ of the undivided energy, down to $U(1)\cdot O(3)$; nothing Lorentzian is broken. $c$ is the unique
conversion between the triples compatible with that residual symmetry and enters only as the weight $c^2$; it is not a symmetry, and $m=E_0/c^2$ is a
rest-mass density of a field point. The reading fails in two respects as first stated: a field point is not a particle (it carries two boost factors,
and the mass of a composite is not interference), and $b$ does not generate $m$ from the undivided energy; at three constituents it moves rest energy
between constituents at fixed totals. Whether some dynamics selects the rest frame or the value of $E_0$ is open; nothing here does.
